\documentclass[10pt,a4paper]{amsart}
\usepackage[margin=19mm]{geometry}
\usepackage{amsmath,amssymb,mathtools}
\usepackage[scr=rsfs,cal=euler]{mathalfa}
\usepackage{xfrac}
\usepackage[unicode,hidelinks,bookmarks=false]{hyperref}
\newcommand{\ie}{i.e.\ }
\newcommand{\cf}{cf.\ }
\newcommand{\resp}{respectively\ }
\newcommand{\Subsection}{\subsection}
\providecommand{\nobreakdash}{\nobreak\hbox{-}}
\setlength{\emergencystretch}{2em}
\multlinegap0pt
\usepackage{bm}
\usepackage{bbm}
\usepackage{dsfont}
\usepackage{xspace}
\usepackage{cancel}
\usepackage{mathtools}
\usepackage{amsthm}
\makeatletter
\@ifundefined{theorem}{\newtheorem{theorem}{Theorem}}{}
\makeatother
\title[Surface-Encoded Partial Coherence Transformation: Modeling S]{Surface-Encoded Partial Coherence Transformation: Modeling Source Coherence Effects in Wave Optics}
\author{Netzer Moriya}
\date{Source: arXiv:2505.17754v1 (CC BY 4.0)}
\begin{document}
\maketitle

\noindent\emph{Source and licence.} Surface-Encoded Partial Coherence Transformation: Modeling Source Coherence Effects in Wave Optics. Version arXiv:2505.17754v1 (CC BY 4.0), distributed under the Creative Commons Attribution 4.0 International licence, as recorded in the arXiv licence metadata for this exact version and shown on the abstract page https://arxiv.org/abs/2505.17754v1. Full rights notice: licenses/arxiv-2505.17754-CC-BY-4.0.txt.\par\smallskip

\noindent\textit{Editorial scope.} Counted unit: the Theorem whose source label is \texttt{None} in the article (source0.tex lines 1429--1440, source label \texttt{None}), delivered together with the article's own complete original proof of it (source0.tex lines 1442--1614, one \texttt{proof} environment of 173 lines). No mathematics is written, repaired or re-derived: statement and proof are the article's own text line for line. Required context retained uncounted: eq:mutual\_coherence (lines 1629--1632). Every definition, display and label of those lines is retained unchanged. Omitted from this package and not redistributed: 1--1428, 1441--1441, 1615--1628, 1633--1774. Nothing inside a retained excerpt is omitted or altered: every retained body is delivered exactly as the article's lines are written, so no omission receipt of any kind accompanies this package. Citations: the retained excerpts carry no citation command at all, so no bibliography entry of the article is transcribed and no reference is redistributed. Reference adaptation: none was needed and none was made; every reference inside the delivered unit resolves inside it, so no unresolved reference or citation is produced. Printed slips kept literal: no printed word, symbol, hypothesis or number of the counted statement or of its proof was altered. Layout and class: the article's own class is \texttt{article} with packages including amsmath, amsfonts, amsthm, amssymb, inputenc, fontenc, graphicx, caption, color, hyperref, breakurl, grffile, makeidx, acronym; the wrapper is the lane's pinned \texttt{amsart} editorial wrapper at 10pt on A4, which restates the theorem-like environments the retained text needs and adds the article's own macro definitions that the retained text needs (none), with extra wrapper packages (bm, bbm, dsfont, xspace, cancel, mathtools). The delivered numbering is the unit's own. Assets: no retained span contains a figure environment, a graphics-inclusion command, a PDF-page inclusion, a table or a TikZ picture; no image file is copied into this package, the package declares no asset dependency and no image file is redistributed with it. Licence: version arXiv:2505.17754v1 of this article is distributed under the Creative Commons Attribution 4.0 International licence, as recorded in the arXiv licence metadata for this exact version and shown on the abstract page https://arxiv.org/abs/2505.17754v1. A full rights notice is delivered at licenses/arxiv-2505.17754-CC-BY-4.0.txt. Characters: every retained excerpt is pure ASCII, so the delivered engine, XeLaTeX with its default Latin Modern text font, reports no missing character.
\begin{equation}
\Gamma_0(\bm{r}_1, \bm{r}_2) = \mathrm{FT}\left[ I_s(\bm{\rho}) \right] \left( \frac{\bm{r}_2 - \bm{r}_1}{\lambda z} \right),
\label{eq:mutual_coherence}
\end{equation}

\begin{theorem}[SECT-VCZ Equivalence with Detailed Steps]
Given a fully coherent incident field $U_i(\mathbf{r}) = A_0$ (uniform plane wave illumination) impinging on a 
surface with coherence operator $\mathcal{C}_S$ characterized by kernel $K_S(\mathbf{r}, \mathbf{r}', \lambda)$, 
followed by propagation operator $\mathcal{P}_z$, the mutual coherence function at the detection plane is mathematically 
equivalent to that obtained from the VCZ theorem for an incoherent source with intensity distribution $I_s(\mathbf{\rho})$, 
provided that:

\begin{equation}
K_S(\mathbf{r}, \mathbf{r}', \lambda) = \mathcal{F}\left\{I_s\left(\frac{\mathbf{\rho}}{\lambda z}\right)\right\}(\mathbf{r} - \mathbf{r}')
\end{equation}
where $\mathcal{F}$ denotes the Fourier transform operator.
\end{theorem}

\begin{proof}
We start with the mutual coherence function at the detection plane using our SECT framework:
\begin{align}
\Gamma_{\text{SECT}}(\mathbf{r}_1, \mathbf{r}_2) &= \langle U_d^*(\mathbf{r}_1) U_d(\mathbf{r}_2) \rangle
\end{align}

Where $U_d(\mathbf{r})$ is the field at the detection plane after applying the reflection operator $\mathcal{R}$, the surface coherence operator $\mathcal{C}_S$, and the propagation operator $\mathcal{P}_z$. We can express this as:
\begin{align}
U_d(\mathbf{r}) &= \mathcal{P}_z(\mathcal{C}_S(\mathcal{R}(U_i)))(\mathbf{r})
\end{align}

With the explicit Fresnel propagator, this becomes:
\begin{align}
U_d(\mathbf{r}) &= \int h(\mathbf{r}, \mathbf{r}'', z, \lambda) \left( \int K_S(\mathbf{r}'', \mathbf{r}', \lambda) \mathcal{R}(U_i)(\mathbf{r}') \, d^2r' \right) \, d^2r'' \\
&= \int \frac{e^{ikz}}{i\lambda z} e^{i\frac{k}{2z}|\mathbf{r} - \mathbf{r}''|^2} \left( \int K_S(\mathbf{r}'', \mathbf{r}', \lambda) \mathcal{R}(U_i)(\mathbf{r}') \, d^2r' \right) \, d^2r'' \notag
\end{align}
where $k = 2\pi/\lambda$ is the wavenumber.

For simplicity, we assume a deterministic reflection operator $\mathcal{R}$ with uniform reflection coefficient $R$ and phase change $\phi_R$:
\begin{align}
\mathcal{R}(U_i)(\mathbf{r}) = R e^{i\phi_R} U_i(\mathbf{r}) = R e^{i\phi_R} A_0 \equiv B_0
\end{align}
where we define $B_0 = R e^{i\phi_R} A_0$ for notational simplicity.

To accurately model the statistical properties of partial coherence, we recognize that the surface coherence operator introduces statistical variations while preserving the mean field. The key insight is that while the average field at the surface may remain $U_s(\mathbf{r}) = B_0$, the surface coherence operator introduces spatial correlations that affect the second-order statistics. We represent this by considering the field at the surface as:
\begin{align}
U_s(\mathbf{r}) = B_0 + \Delta U_s(\mathbf{r})
\end{align}
where $\langle \Delta U_s(\mathbf{r}) \rangle = 0$ and $\langle \Delta U_s^*(\mathbf{r}_1) \Delta U_s(\mathbf{r}_2) \rangle = |B_0|^2 [\gamma_S(\mathbf{r}_1, \mathbf{r}_2) - 1]$.

Here, $\gamma_S(\mathbf{r}_1, \mathbf{r}_2)$ is the normalized mutual coherence function at the surface, which for a spatially stationary kernel depends only on the separation $\mathbf{r}_1 - \mathbf{r}_2$. For the specific kernel form stated in the theorem:
\begin{align}
\gamma_S(\mathbf{r}_1, \mathbf{r}_2) = \mathcal{F}\left\{I_s\left(\frac{\mathbf{\rho}}{\lambda z}\right)\right\}(\mathbf{r}_1 - \mathbf{r}_2)
\end{align}

After propagation to the detection plane using the Fresnel propagator, the field becomes:
\begin{align}
U_d(\mathbf{r}) &= \int \frac{e^{ikz}}{i\lambda z} e^{i\frac{k}{2z}|\mathbf{r} - \mathbf{r}'|^2} [B_0 + \Delta U_s(\mathbf{r}')] \, d^2r' \\
&= B_0 \int \frac{e^{ikz}}{i\lambda z} e^{i\frac{k}{2z}|\mathbf{r} - \mathbf{r}'|^2} \, d^2r' + \int \frac{e^{ikz}}{i\lambda z} e^{i\frac{k}{2z}|\mathbf{r} - \mathbf{r}'|^2} \Delta U_s(\mathbf{r}') \, d^2r' \notag
\end{align}

For the first integral, using the properties of the Fresnel propagator for a uniform field, we have:
\begin{equation}
\int \frac{e^{ikz}}{i\lambda z} e^{i\frac{k}{2z}|\mathbf{r} - \mathbf{r}'|^2} \, d^2r' = e^{ikz}
\end{equation}

Therefore:
\begin{align}
U_d(\mathbf{r}) &= B_0 e^{ikz} + \int \frac{e^{ikz}}{i\lambda z} e^{i\frac{k}{2z}|\mathbf{r} - \mathbf{r}'|^2} \Delta U_s(\mathbf{r}') \, d^2r'
\end{align}

The mutual coherence function at the detection plane is:
\begin{align}
\Gamma_d(\mathbf{r}_1, \mathbf{r}_2) &= \langle U_d^*(\mathbf{r}_1) U_d(\mathbf{r}_2) \rangle \\
&= \langle \left[B_0^* e^{-ikz} + \int \frac{e^{-ikz}}{-i\lambda z} e^{-i\frac{k}{2z}|\mathbf{r}_1 - \mathbf{r}'_1|^2} \Delta U_s^*(\mathbf{r}'_1) \, d^2r'_1 \right] \notag \\
&\quad \times \left[B_0 e^{ikz} + \int \frac{e^{ikz}}{i\lambda z} e^{i\frac{k}{2z}|\mathbf{r}_2 - \mathbf{r}'_2|^2} \Delta U_s(\mathbf{r}'_2) \, d^2r'_2 \right] \rangle \notag
\end{align}

Expanding this product and noting that $\langle \Delta U_s(\mathbf{r}) \rangle = 0$, we get:
\begin{align}
\Gamma_d(\mathbf{r}_1, \mathbf{r}_2) &= |B_0|^2 + \langle \int\int \frac{e^{-ikz}}{-i\lambda z} \frac{e^{ikz}}{i\lambda z} e^{-i\frac{k}{2z}|\mathbf{r}_1 - \mathbf{r}'_1|^2} e^{i\frac{k}{2z}|\mathbf{r}_2 - \mathbf{r}'_2|^2} \\
&\quad \times \Delta U_s^*(\mathbf{r}'_1) \Delta U_s(\mathbf{r}'_2) \, d^2r'_1 \, d^2r'_2 \rangle \notag
\end{align}

Simplifying the pre-factors and substituting the expression for the correlation of $\Delta U_s$:
\begin{align}
\Gamma_d(\mathbf{r}_1, \mathbf{r}_2) &= |B_0|^2 + \frac{1}{(\lambda z)^2} \int\int e^{-i\frac{k}{2z}|\mathbf{r}_1 - \mathbf{r}'_1|^2} e^{i\frac{k}{2z}|\mathbf{r}_2 - \mathbf{r}'_2|^2} \\
&\quad \times |B_0|^2 [\gamma_S(\mathbf{r}'_1, \mathbf{r}'_2) - 1] \, d^2r'_1 \, d^2r'_2 \notag
\end{align}

We can rewrite this as:
\begin{align}
\Gamma_d(\mathbf{r}_1, \mathbf{r}_2) &= |B_0|^2 \left[1 + \frac{1}{(\lambda z)^2} \int\int e^{-i\frac{k}{2z}|\mathbf{r}_1 - \mathbf{r}'_1|^2} e^{i\frac{k}{2z}|\mathbf{r}_2 - \mathbf{r}'_2|^2} [\gamma_S(\mathbf{r}'_1, \mathbf{r}'_2) - 1] \, d^2r'_1 \, d^2r'_2 \right]
\end{align}

For the specific form of $\gamma_S$ given by the Fourier transform of the source intensity distribution:
\begin{align}
\gamma_S(\mathbf{r}'_1, \mathbf{r}'_2) = \mathcal{F}\left\{I_s\left(\frac{\mathbf{\rho}}{\lambda z}\right)\right\}(\mathbf{r}'_1 - \mathbf{r}'_2)
\end{align}

We can apply the convolution theorem to evaluate the double integral. Let's denote:
\begin{align}
f_1(\mathbf{r}') &= e^{-i\frac{k}{2z}|\mathbf{r}_1 - \mathbf{r}'|^2} \\
f_2(\mathbf{r}') &= e^{i\frac{k}{2z}|\mathbf{r}_2 - \mathbf{r}'|^2}
\end{align}

In the paraxial approximation, these can be expanded as:
\begin{align}
f_1(\mathbf{r}') &= e^{-i\frac{k}{2z}(|\mathbf{r}_1|^2 + |\mathbf{r}'|^2 - 2\mathbf{r}_1 \cdot \mathbf{r}')} = e^{-i\frac{k}{2z}|\mathbf{r}_1|^2} e^{-i\frac{k}{2z}|\mathbf{r}'|^2} e^{i\frac{k}{z}\mathbf{r}_1 \cdot \mathbf{r}'} \\
f_2(\mathbf{r}') &= e^{i\frac{k}{2z}(|\mathbf{r}_2|^2 + |\mathbf{r}'|^2 - 2\mathbf{r}_2 \cdot \mathbf{r}')} = e^{i\frac{k}{2z}|\mathbf{r}_2|^2} e^{i\frac{k}{2z}|\mathbf{r}'|^2} e^{-i\frac{k}{z}\mathbf{r}_2 \cdot \mathbf{r}'}
\end{align}

Now let's define the integral:
\begin{align}
I(\mathbf{r}_1, \mathbf{r}_2) &= \int\int f_1(\mathbf{r}'_1) f_2(\mathbf{r}'_2) \gamma_S(\mathbf{r}'_1, \mathbf{r}'_2) \, d^2r'_1 \, d^2r'_2 \\
&= \int\int e^{-i\frac{k}{2z}|\mathbf{r}_1|^2} e^{-i\frac{k}{2z}|\mathbf{r}'_1|^2} e^{i\frac{k}{z}\mathbf{r}_1 \cdot \mathbf{r}'_1} e^{i\frac{k}{2z}|\mathbf{r}_2|^2} e^{i\frac{k}{2z}|\mathbf{r}'_2|^2} e^{-i\frac{k}{z}\mathbf{r}_2 \cdot \mathbf{r}'_2} \notag \\
&\quad \times \mathcal{F}\left\{I_s\left(\frac{\mathbf{\rho}}{\lambda z}\right)\right\}(\mathbf{r}'_1 - \mathbf{r}'_2) \, d^2r'_1 \, d^2r'_2 \notag
\end{align}

Factoring out terms that don't depend on the integration variables:
\begin{align}
I(\mathbf{r}_1, \mathbf{r}_2) &= e^{-i\frac{k}{2z}|\mathbf{r}_1|^2} e^{i\frac{k}{2z}|\mathbf{r}_2|^2} \int\int e^{-i\frac{k}{2z}|\mathbf{r}'_1|^2} e^{i\frac{k}{z}\mathbf{r}_1 \cdot \mathbf{r}'_1} e^{i\frac{k}{2z}|\mathbf{r}'_2|^2} e^{-i\frac{k}{z}\mathbf{r}_2 \cdot \mathbf{r}'_2} \\
&\quad \times \mathcal{F}\left\{I_s\left(\frac{\mathbf{\rho}}{\lambda z}\right)\right\}(\mathbf{r}'_1 - \mathbf{r}'_2) \, d^2r'_1 \, d^2r'_2 \notag
\end{align}

Making the change of variables $\mathbf{u} = \mathbf{r}'_1$ and $\mathbf{v} = \mathbf{r}'_1 - \mathbf{r}'_2$, with Jacobian $d^2r'_1 \, d^2r'_2 = d^2u \, d^2v$:
\begin{align}
I(\mathbf{r}_1, \mathbf{r}_2) &= e^{-i\frac{k}{2z}|\mathbf{r}_1|^2} e^{i\frac{k}{2z}|\mathbf{r}_2|^2} \int\int e^{-i\frac{k}{2z}|\mathbf{u}|^2} e^{i\frac{k}{z}\mathbf{r}_1 \cdot \mathbf{u}} e^{i\frac{k}{2z}|\mathbf{u}-\mathbf{v}|^2} e^{-i\frac{k}{z}\mathbf{r}_2 \cdot (\mathbf{u}-\mathbf{v})} \\
&\quad \times \mathcal{F}\left\{I_s\left(\frac{\mathbf{\rho}}{\lambda z}\right)\right\}(\mathbf{v}) \, d^2u \, d^2v \notag
\end{align}

Expanding $|\mathbf{u}-\mathbf{v}|^2 = |\mathbf{u}|^2 + |\mathbf{v}|^2 - 2\mathbf{u} \cdot \mathbf{v}$:
\begin{align}
I(\mathbf{r}_1, \mathbf{r}_2) &= e^{-i\frac{k}{2z}|\mathbf{r}_1|^2} e^{i\frac{k}{2z}|\mathbf{r}_2|^2} \int\int e^{-i\frac{k}{2z}|\mathbf{u}|^2} e^{i\frac{k}{z}\mathbf{r}_1 \cdot \mathbf{u}} e^{i\frac{k}{2z}|\mathbf{u}|^2} e^{i\frac{k}{2z}|\mathbf{v}|^2} e^{-i\frac{k}{z}\mathbf{u} \cdot \mathbf{v}} \\
&\quad \times e^{-i\frac{k}{z}\mathbf{r}_2 \cdot \mathbf{u}} e^{i\frac{k}{z}\mathbf{r}_2 \cdot \mathbf{v}} \mathcal{F}\left\{I_s\left(\frac{\mathbf{\rho}}{\lambda z}\right)\right\}(\mathbf{v}) \, d^2u \, d^2v \notag
\end{align}

Simplifying and rearranging terms:
\begin{align}
I(\mathbf{r}_1, \mathbf{r}_2) &= e^{-i\frac{k}{2z}|\mathbf{r}_1|^2} e^{i\frac{k}{2z}|\mathbf{r}_2|^2} \int e^{i\frac{k}{2z}|\mathbf{v}|^2} e^{i\frac{k}{z}\mathbf{r}_2 \cdot \mathbf{v}} \mathcal{F}\left\{I_s\left(\frac{\mathbf{\rho}}{\lambda z}\right)\right\}(\mathbf{v}) \\
&\quad \times \left[\int e^{i\frac{k}{z}(\mathbf{r}_1-\mathbf{r}_2-\mathbf{v}) \cdot \mathbf{u}} \, d^2u \right] \, d^2v \notag
\end{align}

The inner integral evaluates to $(2\pi)^2 \delta(\frac{k}{z}(\mathbf{r}_1-\mathbf{r}_2-\mathbf{v}))$, where $\delta$ is 
the Dirac delta function. This gives:
\begin{align}
I(\mathbf{r}_1, \mathbf{r}_2) &= e^{-i\frac{k}{2z}|\mathbf{r}_1|^2} e^{i\frac{k}{2z}|\mathbf{r}_2|^2} \left(\frac{z}{k}\right)^2 (2\pi)^2 \\
&\quad \times \int e^{i\frac{k}{2z}|\mathbf{v}|^2} e^{i\frac{k}{z}\mathbf{r}_2 \cdot \mathbf{v}} \mathcal{F}\left\{I_s\left(\frac{\mathbf{\rho}}{\lambda z}\right)\right\}(\mathbf{v}) \delta(\mathbf{r}_1-\mathbf{r}_2-\mathbf{v}) \, d^2v \notag
\end{align}

Evaluating the remaining integral using the sifting property of the delta function:
\begin{align}
I(\mathbf{r}_1, \mathbf{r}_2) &= e^{-i\frac{k}{2z}|\mathbf{r}_1|^2} e^{i\frac{k}{2z}|\mathbf{r}_2|^2} \left(\frac{z}{k}\right)^2 (2\pi)^2 \\
&\quad \times e^{i\frac{k}{2z}|\mathbf{r}_1-\mathbf{r}_2|^2} e^{i\frac{k}{z}\mathbf{r}_2 \cdot (\mathbf{r}_1-\mathbf{r}_2)} \mathcal{F}\left\{I_s\left(\frac{\mathbf{\rho}}{\lambda z}\right)\right\}(\mathbf{r}_1-\mathbf{r}_2) \notag
\end{align}

Further simplifying using $|\mathbf{r}_1|^2 = |\mathbf{r}_2|^2 + |\mathbf{r}_1-\mathbf{r}_2|^2 + 2\mathbf{r}_2 \cdot (\mathbf{r}_1-\mathbf{r}_2)$:
\begin{align}
I(\mathbf{r}_1, \mathbf{r}_2) &= \left(\frac{z}{k}\right)^2 (2\pi)^2 \mathcal{F}\left\{I_s\left(\frac{\mathbf{\rho}}{\lambda z}\right)\right\}(\mathbf{r}_1-\mathbf{r}_2)
\end{align}

Therefore, returning to the expression for the mutual coherence function~\ref{eq:mutual_coherence}:
\begin{align}
\Gamma_d(\mathbf{r}_1, \mathbf{r}_2) &= |B_0|^2 \left[1 + \frac{1}{(\lambda z)^2} \left(I(\mathbf{r}_1, \mathbf{r}_2) - I_0\right)\right] \\
&= |B_0|^2 \left[1 + \frac{1}{(\lambda z)^2} \left(\left(\frac{z}{k}\right)^2 (2\pi)^2 \mathcal{F}\left\{I_s\left(\frac{\mathbf{\rho}}{\lambda z}\right)\right\}(\mathbf{r}_1-\mathbf{r}_2) - I_0\right)\right]
\end{align}

where $I_0$ is the corresponding integral for the constant term ($\gamma_S = 1$).

Substituting $k = 2\pi/\lambda$ and simplifying:
\begin{align}
\Gamma_d(\mathbf{r}_1, \mathbf{r}_2) &= |B_0|^2 \left[1 + \frac{1}{(\lambda z)^2} \left(\frac{z^2 \lambda^2}{4\pi^2} (2\pi)^2 \mathcal{F}\left\{I_s\left(\frac{\mathbf{\rho}}{\lambda z}\right)\right\}(\mathbf{r}_1-\mathbf{r}_2) - I_0\right)\right] \\
&= |B_0|^2 \left[1 + \mathcal{F}\left\{I_s\left(\frac{\mathbf{\rho}}{\lambda z}\right)\right\}(\mathbf{r}_1-\mathbf{r}_2) - \frac{I_0}{(\lambda z)^2}\right]
\end{align}

The term $I_0$ corresponds to the normalization constant, and after proper normalization:
\begin{align}
\Gamma_d(\mathbf{r}_1, \mathbf{r}_2) &= |B_0|^2 \mathcal{F}\left\{I_s\left(\frac{\mathbf{\rho}}{\lambda z}\right)\right\}(\mathbf{r}_1-\mathbf{r}_2)
\end{align}

This is precisely the form predicted by the Van Cittert-Zernike theorem for an incoherent source with intensity distribution $I_s(\mathbf{\rho})$:
\begin{align}
\Gamma_{\text{VCZ}}(\mathbf{r}_1, \mathbf{r}_2) &= C \mathcal{F}\left\{I_s\left(\frac{\mathbf{\rho}}{\lambda z}\right)\right\}(\mathbf{r}_1-\mathbf{r}_2)
\end{align}
where $C$ is a normalization constant.

Therefore:
\begin{align}
\Gamma_{\text{SECT}}(\mathbf{r}_1, \mathbf{r}_2) \propto \Gamma_{\text{VCZ}}(\mathbf{r}_1, \mathbf{r}_2)
\end{align}

This establishes that our dual-component SECT framework can exactly reproduce the coherence properties predicted by the VCZ theorem when the appropriate surface coherence kernel is used.
\end{proof}

\end{document}
